% =====================================================================
% Capítulo 2 (continuación): Secs. 2.9-2.10
%   Estado del arte (trabajos desde 2020) y consideraciones finales
% =====================================================================

% ---------------------------------------------------------------------
\section{Estado del arte}
\label{sec:estado-arte}

Esta sección revisa los trabajos publicados desde 2020 que se relacionan con la
propuesta. No incluye a \ac{ST-CC} ni a SigTree, que se presentaron en las
Secciones~\ref{sec:stcc} y~\ref{sec:sigtree} como los dos métodos de los que la
propuesta parte directamente: conserva la partición base de \ac{ST-CC} y le
añade un refinamiento local cuyo criterio de parada, inspirado en SigTree, es una
prueba de hipótesis y no un umbral heurístico. Los trabajos revisados se
organizan en tres familias: el co-agrupamiento espacio-temporal y las
estructuras de bloques flexibles, la inferencia selectiva para agrupamientos y
estructuras de bloques, y la detección de focos de actividad (\emph{hotspots})
espacio-temporales. Al final de cada familia se indica en qué se diferencia de la
propuesta, y la sección cierra con una síntesis comparativa que identifica la
brecha de investigación.

% ---------------------------------------------------------------------
\subsection{Co-agrupamiento espacio-temporal y estructuras de bloques flexibles}
\label{sec:ea_coagrupamiento}

La primera familia agrupa dos tipos de trabajos que atacan el mismo
problema de dos maneras distintas: aplicar co-agrupamiento o biclustering a
datos con una dimensión espacial y otra temporal, o bien relajar la rejilla
regular decidiendo de forma flexible cuántos bloques usar.

Dentro del mismo marco de Bregman en el que se basa \ac{ST-CC}
(Sección~\ref{sec:stcc}), \cite{Wu2022} propuso BCC\_MSSR, una
variante que usa el residuo cuadrático para encontrar co-clusters con
\emph{tendencias coherentes} en lugar de valores similares, y la aplicó a
temperaturas mensuales de estaciones meteorológicas. Este trabajo muestra
que el marco de Bregman admite distintas nociones de homogeneidad dentro de
un bloque, algo relevante para esta tesis, que define la heterogeneidad de
un bloque mediante la interacción residual entre filas y columnas.

En movilidad vial, \cite{Neves2022} utilizaron biclustering para extraer
patrones accionables a partir de datos de tráfico heterogéneos. A diferencia
del co-agrupamiento en tablero, el biclustering busca submatrices que no
necesitan cubrir toda la matriz. Cabe destacar que este trabajo sí
incorpora una prueba de significancia estadística (BSig) para retener solo
biclusters cuyo patrón de congestión sea improbable bajo un modelo nulo,
con un nivel de significancia de 0.01 por bicluster.

Un segundo grupo de trabajos, en lugar de cambiar la noción de coherencia
dentro del bloque, flexibiliza directamente el número o la forma de la
estructura de bloques. En redes bipartitas, que equivalen a una matriz de
filas y columnas, \cite{Yen2020} propusieron una formulación bayesiana no
paramétrica del modelo de bloques estocástico (biSBM) que elige de forma
parsimoniosa el número de comunidades. Al compararla con un modelo de
bloques \emph{jerárquico} de alta resolución, encontraron un régimen,
formado por redes pequeñas y dispersas, en el que los modelos no
jerárquicos superan a su contraparte más flexible. Este hallazgo advierte
que más flexibilidad estructural no garantiza mejores resultados cuando los
datos son escasos.

En cuanto a algoritmos jerárquicos, \cite{Sugahara2023} propusieron HICCAM,
un método de co-agrupamiento jerárquico que incorpora matrices aumentadas
con información de dominios externos. Con un enfoque frecuentista,
\cite{Watanabe2021gof} desarrollaron una prueba de bondad de ajuste para
modelos de bloques latentes que evalúa si la matriz observada se ajusta a un
número dado de grupos de filas y columnas, o si contiene más grupos en al
menos una dirección. La prueba se apoya en resultados de la teoría de
matrices aleatorias para la matriz de covarianza muestral.

\textbf{Diferencia con la propuesta.} Estos cinco trabajos cubren dos
maneras de escapar de la rejilla fija, pero ninguna coincide con la de esta
tesis. Por un lado, los métodos basados en Bregman producen particiones
globales en tablero, y el biclustering de \cite{Neves2022}, aunque sí
incorpora una prueba de significancia, la aplica como filtro
\emph{posterior} sobre un conjunto de candidatos generado de forma
exhaustiva a distintos niveles de soporte y coherencia, no como un criterio
que decide, durante la propia búsqueda, si una submatriz concreta merece
subdividirse. Por otro lado, los métodos que flexibilizan el número de
bloques lo hacen por vía bayesiana, mediante una jerarquía construida sobre
la totalidad de filas y columnas, o eligiendo un número \emph{global} de
grupos; en ningún caso el resultado refina un bloque \emph{desde dentro}.
Ninguno de los cinco trabajos usa una prueba estadística como criterio de
\emph{parada} que gobierne, bloque por bloque, la construcción de la
partición. La propuesta conserva el marco de \ac{ST-CC} y le añade un
refinamiento local en el que la prueba de hipótesis gobierna la
construcción misma del árbol: la jerarquía de columnas queda anidada dentro
de un subconjunto de filas, por lo que el resultado deja de ser una
rejilla. Por ello, un experimento necesario en esta tesis es comparar el
refinamiento local contra una rejilla global más fina con el mismo número
de parámetros.

% ---------------------------------------------------------------------
\subsection{Inferencia selectiva para agrupamientos y estructuras de bloques}
\label{sec:ea_inferencia}

Mientras que SigTree (Sección~\ref{sec:sigtree}) usa pruebas de
hipótesis para \emph{construir} una partición, esta segunda familia las usa
para \emph{validar} agrupamientos o estructuras de bloques que ya fueron
estimados a partir de los mismos datos con los que después se prueban.

\cite{Gao2024} mostraron que, cuando los grupos se definen mediante
agrupamiento, una prueba clásica de diferencia de medias produce una tasa de
error tipo~I extremadamente inflada, y que el problema persiste incluso si
se usan dos conjuntos de datos independientes para definir y para probar los
grupos. Para resolverlo, propusieron una prueba selectiva que calcula
$p$-valores exactos para grupos obtenidos por agrupamiento jerárquico
aglomerativo. \cite{Chen2023} extendieron la idea a \textit{K-means}, con un
$p$-valor que condiciona sobre todas las asignaciones intermedias del
algoritmo.

Este problema de doble uso de los datos también aparece en matrices con
estructura de bloques y en datos de conteo. \cite{Watanabe2021si}
propusieron un método de inferencia selectiva para modelos de bloques
latentes. Su prueba evalúa las pertenencias de filas y columnas obtenidas
por minimización de residuos cuadráticos, sirve también para probar el
número de grupos, e incluye una versión aproximada basada en recocido
simulado para evitar la explosión combinatoria.

Una alternativa más sencilla es dividir cada observación en lugar de
dividir la muestra. \cite{Neufeld2024cs} introdujeron el \emph{count
splitting}, que permite estimar una variable latente (por ejemplo, un
agrupamiento) con una parte de los conteos y probarla con otra, bajo un
supuesto de Poisson. \cite{Neufeld2024dt} generalizaron esta idea con el
\emph{data thinning}, aplicable a cualquier distribución cerrada bajo
convolución, como la Poisson o la binomial negativa, y mostraron que sirve
para validar resultados de aprendizaje no supervisado.

\textbf{Diferencia con la propuesta.} Estos cinco trabajos validan una
estructura ya estimada, en una sola dimensión o para una matriz de bloques
completa, pero ninguno construye ni valida un refinamiento jerárquico
local. La propuesta traslada el principio de SigTree —construir la
partición con una prueba de hipótesis, no solo validarla después— a
bloques de una matriz de conteos, donde la división puede ser de filas, de
columnas o conjunta. Se apoya en la prueba de \cite{Watanabe2021si} como la
referencia más rigurosa para estructuras de bloques —aunque planteada para
residuos cuadráticos y no para conteos— y en la idea de separar la selección de
la inferencia del \emph{count splitting} de \cite{Neufeld2024cs}, que adapta a
matrices de conteos de viajes dividiendo los días en lugar de los conteos,
porque estos presentan sobredispersión. El sesgo de selección demostrado por \cite{Gao2024}
exige además repetir la búsqueda del mejor corte dentro de cada
permutación. Los fundamentos de esta línea se presentan en la
Sección~\ref{sec:est-seleccion}.

% ---------------------------------------------------------------------
\subsection{Detección de hotspots espacio-temporales}
\label{sec:ea_hotspots}

La última familia aborda la detección de concentraciones de eventos en el
espacio y el tiempo, que es el tipo de patrón local que esta tesis busca
revelar. \cite{Shiode2022} propusieron \emph{space-time NetScan}, una
estadística de barrido (\emph{scan statistic}) que mide los patrones
espacio-temporales a lo largo de la red de calles. En pruebas con datos
sintéticos, el método superó a los existentes al detectar la ubicación,
forma, tamaño y duración de los focos. Los autores también señalan que la
corrección de Bonferroni usada en trabajos previos resultaba demasiado
conservadora para detectar focos.
% NOTA: esta última afirmación (Bonferroni "demasiado conservadora") no
% pude verificarla en el fragmento del paper que revisé (solo las
% primeras 2 páginas); es plausible y consistente con la literatura de
% scan statistics, pero confirma la página exacta antes de citarla así
% de categóricamente.

En el dominio de los taxis, \cite{Chen2021hotspots} abordaron la detección
de focos urbanos a partir de las paradas de taxi mediante un criterio de
densidad máxima local. El umbral de popularidad de cada hotspot se fija
mediante una simulación Monte Carlo comparada con un particionamiento de
head/tail breaks, un procedimiento basado en aleatorización pero sin un
nivel de significancia $\alpha$ fijo ni corrección por comparaciones
múltiples. \cite{Mishra2025}, por su parte, observaron que los métodos que
solo agrupan zonas de alta demanda concentran los focos en el centro de la
ciudad e ignoran el resto. Para corregirlo, formularon la detección como un
problema de conjunto dominante de $k$ saltos en el espacio espacio-temporal,
cubriendo también las zonas de demanda moderada en distintas horas del día,
y evaluaron su método con datos de taxis de Nueva York.
% NOTA: verificar año y venue exactos de esta cita. Existe una línea de
% trabajos muy similares del mismo autor (detección de hotspots de taxi
% vía dominating sets) publicados en distintos años; no logré confirmar
% con certeza cuál corresponde específicamente a 2025.

\textbf{Diferencia con la propuesta.} Estos métodos buscan zonas de alta
intensidad en un espacio continuo, en una red de calles o en un grafo,
mientras que la propuesta trabaja sobre una matriz región~$\times$~tiempo y
detecta interacciones residuales dentro de bloques ya formados. Las
estadísticas de barrido comparten con la propuesta la lógica de la razón de
verosimilitud y la necesidad de controlar las comparaciones múltiples, por
lo que son un punto de comparación natural. El trabajo de
\cite{Mishra2025}, además, confirma que los focos de interés no siempre
están donde la demanda es mayor, lo que refuerza la necesidad de un
criterio basado en la interacción y no solo en el volumen.

% ---------------------------------------------------------------------
\subsection{Síntesis y brecha de investigación}
\label{sec:ea_sintesis}

La Tabla~\ref{tab:comparacion_ea} resume las características de los
trabajos de las tres familias de literatura relacionada; no incluye a
\ac{ST-CC} ni a SigTree, presentados en las Secciones~\ref{sec:stcc}
y~\ref{sec:sigtree} como los métodos de los que la propuesta parte directamente, en vez de
como literatura de la que se diferencia. Las columnas indican si el
trabajo trata datos espacio-temporales, si particiona filas y columnas a la
vez, si permite estructuras locales o jerárquicas, si usa un criterio
estadístico para decidir o validar la estructura, y si controla el error
ante múltiples pruebas o ante el sesgo de selección.

\begin{table}[htbp]
\centering
\small
\setlength{\tabcolsep}{4pt}
\begin{tabular}{>{\raggedright\arraybackslash}p{4.2cm} *{5}{>{\centering\arraybackslash}p{1.9cm}}}
\hline
\textbf{Trabajo} & \textbf{Datos esp.-temp.} & \textbf{Filas y columnas}
& \textbf{Estructura local} & \textbf{Criterio estadístico}
& \textbf{Control de error} \\
\hline
\multicolumn{6}{l}{\textit{Co-agrupamiento espacio-temporal y estructuras de bloques flexibles}} \\
\cite{Wu2022}                  & $\checkmark$ & $\checkmark$ & $\times$     & $\times$      & $\times$ \\
\cite{Neves2022}               & $\checkmark$ & $\checkmark$ & $\checkmark$ & $\checkmark$ & $\times$ \\
\cite{Yen2020}                 & $\times$     & $\checkmark$ & $\times$     & Bayesiano     & $\times$ \\
\cite{Sugahara2023}            & $\times$     & $\checkmark$ & $\checkmark$ & $\times$ & $\times$ \\
% NOTA: última celda ("Control de error") pendiente de verificación
% directa en la Sección 3 del paper (Watanabe & Suzuki, 2021, CSDA 154):
% no queda claro por el resumen si el procedimiento secuencial sobre
% (K,H) corrige el nivel alfa a lo largo de la secuencia de pruebas.
\cite{Watanabe2021gof}         & $\times$     & $\checkmark$ & $\times$     & $\checkmark$  & Por verificar \\
\multicolumn{6}{l}{\textit{Inferencia selectiva para agrupamientos y estructuras de bloques}} \\
\cite{Gao2024}                 & $\times$     & $\times$     & $\checkmark$ & $\checkmark$  & $\checkmark$ \\
\cite{Chen2023}                & $\times$     & $\times$     & $\times$     & $\checkmark$  & $\checkmark$ \\
\cite{Watanabe2021si}          & $\times$     & $\checkmark$ & $\times$     & $\checkmark$  & $\checkmark$ \\
\cite{Neufeld2024cs}           & $\times$     & $\times$     & $\times$     & $\checkmark$  & $\checkmark$ \\
\cite{Neufeld2024dt}           & $\times$     & $\times$     & $\times$     & $\checkmark$  & $\checkmark$ \\
\multicolumn{6}{l}{\textit{Detección de hotspots espacio-temporales}} \\
\cite{Shiode2022}              & $\checkmark$ & $\times$     & $\checkmark$ & $\checkmark$  & $\checkmark$ \\
\cite{Chen2021hotspots}        & $\checkmark$ & $\times$     & $\checkmark$ & $\times$ & $\times$ \\
\cite{Mishra2025}              & $\checkmark$ & $\times$     & $\checkmark$ & $\times$      & $\times$ \\
\hline
\textbf{Propuesta}             & $\checkmark$ & $\checkmark$ & $\checkmark$ & $\checkmark$  & $\checkmark$ \\
\hline
\end{tabular}
\caption{Comparación de los trabajos revisados con la propuesta. No incluye
a \ac{ST-CC} ni a SigTree, presentados en las Secciones~\ref{sec:stcc}
y~\ref{sec:sigtree}.}
\label{tab:comparacion_ea}
\end{table}

La tabla muestra que las piezas necesarias existen por separado, más allá
de los dos métodos base ya presentados. El co-agrupamiento espacio-temporal
está consolidado, pero produce rejillas globales o estructuras locales sin
criterio de parada estadístico. Las pruebas de hipótesis y la inferencia
post-selección ofrecen criterios rigurosos, pero en una sola dimensión o
para validar estructuras ya formadas. Los métodos de detección de
\emph{hotspots} encuentran focos locales, pero no operan sobre una
partición de co-agrupamiento.

En la revisión realizada no se encontró un trabajo que combine las tres
características: refinar localmente los bloques de un co-agrupamiento
espacio-temporal, decidir cada división mediante una prueba de hipótesis y
controlar el error a lo largo de la jerarquía. Esta es la brecha que aborda
la presente tesis.

% ---------------------------------------------------------------------
\section{Consideraciones Finales}
\label{sec:marco-finales}

Este capítulo reunió los fundamentos de la propuesta y la revisión de los
trabajos relacionados. En la parte teórica se mostró que, con la I-divergencia y
datos de conteo, las bases C2, C5 y C6 del co-agrupamiento de Bregman
corresponden a modelos log-lineales de Poisson anidados, y que la divergencia de
una partición en tablero se descompone en tres términos
(Proposición~\ref{prop:descomposicion}). El último de ellos mide la interacción
entre regiones y horarios que queda dentro de cada bloque; es el objeto que el
método propuesto intenta detectar, y el estadístico $G^2$ que lo mide es la pieza
que conecta el marco de \ac{ST-CC} con las pruebas de hipótesis. La parte
estadística presentó las herramientas para usar ese estadístico como criterio de
decisión: pruebas condicionales por simulación que no dependen de aproximaciones
asintóticas, procedimientos jerárquicos para controlar el error a lo largo de un
árbol, formas de corregir el doble uso de los datos y medidas de tamaño de efecto
que no crecen con la muestra. SigTree mostró que el principio de construir una
partición con pruebas de hipótesis es viable, y sus limitaciones indicaron qué
debe adaptarse.

La revisión del estado del arte organizó los trabajos recientes en tres familias
y mostró que ninguno reúne las características de la propuesta: refinar
localmente los bloques de un co-agrupamiento espacio-temporal, decidir cada
división con una prueba de hipótesis y controlar el error a lo largo de la
jerarquía. Dos lecciones de la revisión orientan además el diseño del método: la
inferencia post-selección obliga a corregir el sesgo que surge al probar
divisiones elegidas con los mismos datos, y la inflación de los estadísticos con
muestras grandes exige complementar la significancia con un umbral de tamaño de
efecto. El Capítulo~\ref{chap:propuesta} presenta el método que integra estos
elementos.